\documentclass[12pt, a4paper]{article}
\usepackage[french]{babel}
\usepackage{helvet}
\usepackage[T1]{fontenc}
\usepackage[utf8]{inputenc}
\usepackage{geometry}
\usepackage{fancyhdr}
\usepackage{titlesec}
\usepackage{array}
\usepackage{graphicx}
\usepackage{amsmath}
\usepackage{xcolor}
\usepackage{enumitem}
\usepackage{hyperref}
\usepackage{booktabs}
\usepackage{tikz}
\usetikzlibrary{calc}
\usepackage{eso-pic}

\geometry{margin=2.5cm}
\titleformat{\section}{\large\bfseries}{\thesection}{1em}{}

\pagestyle{fancy}
\fancyhf{}
\fancyhead[L]{\footnotesize\textbf{Rapport du LAB N°1}}
\fancyhead[R]{\footnotesize\textbf{Théories et Pratiques de l'Investigation Numérique}}
\fancyfoot[C]{\thepage}

\newcommand{\bordurepage}{
    \begin{tikzpicture}[remember picture, overlay]
    \draw[line width=4pt, rounded corners=10pt, color=blue!60] 
        ($(current page.north west) + (1.5cm,-1.5cm)$) 
        rectangle 
        ($(current page.south east) + (-1.5cm,1.5cm)$);
    \end{tikzpicture}
}

\begin{document}

% Page de garde
\begin{titlepage}
	\bordurepage
	\begin{tabular}{p{0.38\linewidth} p{0.2\linewidth} p{0.38\linewidth}}
		\centering
		\begin{footnotesize}
			\textbf{RÉPUBLIQUE DU CAMEROUN}                   \\
			******                                            \\
			Paix \textendash{} Travail \textendash{} Patrie   \\
			******                                            \\
			\textbf{UNIVERSITÉ DE YAOUNDÉ I}                  \\
			******                                            \\
			ÉCOLE NATIONALE SUPÉRIEURE                        \\
			POLYTECHNIQUE DE YAOUNDÉ                          \\
			******                                            \\
			\textbf{DEPARTEMENT DE GENIE INFORMATIQUE}        \\
			******
		\end{footnotesize}
		 &
		\centering
		\vspace{0.5cm}
		\includegraphics[width = 1 \linewidth]{Logo/Logo_ENSPY.png}
		 &
		\centering
		\begin{footnotesize}
			\textbf{REPUBLIC OF CAMEROON}                     \\
			******                                            \\
			Peace \textendash{} Work \textendash{} Fatherland \\
			******                                            \\
			\textbf{UNIVERSITY OF YAOUNDÉ I}                  \\
			******                                            \\
			NATIONAL ADVANCED SCHOOL                          \\
			ENGINEERING OF YAOUNDE                            \\
			******                                            \\
			\textbf {COMPUTERS ENGINEERING DEPARTMENT}        \\
			******
		\end{footnotesize} \\
	\end{tabular}
	\centering
	\vspace*{1cm}

	% Titre principal
	\noindent
	\begin{tikzpicture}
		\node[
			draw=blue!60,
			line width=4pt,
			rounded corners=10pt,
			inner sep=14pt,
			text width=0.9\textwidth,
			align=center
		]{
			{\Huge\bfseries Rapport du LAB N°1}
		};
	\end{tikzpicture}

	\vspace{1cm}

	% Sous-titre
	{\Large \textbf{Configuration d'un Environnement Réseau Fonctionnel et Sécurisé}}

	\vspace{0.5cm}

	% Titre du cours
	{\LARGE \textbf{Théories et Pratiques de l'Investigation Numérique}}

	\vspace{2cm}

	% Section auteurs
	{\large Rédigé par :}

	\vspace{0.5cm}

	\begin{tabular}{|>{\centering\arraybackslash}m{8cm}
		|>{\centering\arraybackslash}m{4cm}
		|>{\centering\arraybackslash}m{3cm}|}
		\hline
		\textbf{Noms \& Prénoms} & \textbf{Filière} & \textbf{Matricule} \\
		\hline
		DSAMAGO JAFFO Trésor     & HN -- CIN 4      & 22P036             \\
		\hline
	\end{tabular}

	\vfill

	% Encadrement
	\begin{Large}
		\textbf{Sous la direction de :} \\
		M. \textbf{MINKA MI NGUIDJOI Thierry Emmanuel} \\
	\end{Large}

	\vspace{1cm}

	% Année académique
	\textbf{Année académique : 2025--2026}
\end{titlepage}

\section{Objectifs du LAB}
Le présent LAB vise à donner aux étudiants l’occasion de configurer une
infrastructure réseau constituée d’un équipement de frontière, d’un réseau
local et d’une DMZ. Les étudiants réviseront ainsi la configuration et le
paramétrage d’équipements actifs, l’adressage réseau et le partage de
ressources. Ce LAB s’inscrit dans un ensemble plus grand visant à réaliser
l’investigation numérique dans une société après qu’elle a été victime d’un
ransomware.

\section{Préparation de l'infrastructure}
Avant de créer et configurer notre infrastructure, nous avions commencé par
créer les machines virtuelles, installer les images des équipements réseaux
(Switch, Router, Parefeu) sur notre logiciel de simulation réseau : GNS3.
\subsection{Création des machines virtuelles}
Pour ce LAB, nous avons créés 03 machines virtuelles à savoir : \textbf{Windows
	10} et \textbf{Windows 11}, qui serviront de machines clientes, \textbf{Ubuntu}
et \textbf{Kali} qui serviront de serveurs à l'aide du Logiciel VM Workstation
pro v17.52.
\subsubsection{Machine virtuelle Windows}
\begin{itemize}
	\item \textbf{Nom de la machine :} Win 10
	\item \textbf{Système d'exploitation :} Windows 10 Professional 22H2
	\item \textbf{RAM :} 4 Go
	\item \textbf{Processeur :} 2 cœurs
	\item \textbf{Disque dur :} 60 Go
\end{itemize}
Pour la machine virtuelle WIndows 11, nous avons simplement cloner la machine virtuelle Windows 10, puis nous avons modifié le nom de la machine en \textbf{Win 11}.
\subsubsection{Serveur Ubuntu}
\begin{itemize}
	\item \textbf{Nom de la machine :} Ubuntu
	\item \textbf{Système d'exploitation :} Ubuntu 22.04 LTS
	\item \textbf{RAM :} 4 Go
	\item \textbf{Processeur :} 2 cœurs
	\item \textbf{Disque dur :} 25 Go
\end{itemize}
Nous avons installé une application Web appélé \textbf{APPLICATION} sur le serveur Ubuntu, qui nous permettra de tester la connectivité entre les machines clientes et le serveur.
\subsubsection{Machine virtuelle Kali}
\begin{itemize}
	\item \textbf{Nom de la machine :} Kali
	\item \textbf{Système d'exploitation :} Kali Linux 2023.4
	\item \textbf{RAM :} 4 Go
	\item \textbf{Processeur :} 2 cœurs
	\item \textbf{Disque dur :} 30 Go
\end{itemize}
\subsection{Equipements d'interconnexions}
Pour l'interconnexion de nos machines virtuelles, nous avons utilisé les
équipements suivants :
\begin{itemize}
	\item Un \textbf{Switch manageable CISCO 2960} pour la connexion de la machines
	      cliente Win 10 et notre Parefeu.
	\item Un \textbf{Routeur CISCO C7200} pour la connexion entre la DMZ et le réseau
	      local (LAN).
	\item Un \textbf{Parefeu FORTINET} pour la sécurisation de l'infrastructure réseau.
\end{itemize}
\section{Réalisation de l'infrastructure}
\subsection{Table d'adressage réseau}
La table d'adressage réseau est un tableau qui permet de définir les adresses
IP des différentes machines de notre infrastructure. Voici la table d'adressage
que nous avons utilisée :
\begin{table}[h!]
	\centering
	\renewcommand{\arraystretch}{1.3} % augmente la hauteur des lignes
	\caption{Plan d'adressage réseau}
	\begin{tabular}{|l|c|c|c|c|}
		\hline
		\textbf{Machine} & \textbf{Adresse IP} & \textbf{Masque} & \textbf{Passerelle} & \textbf{Type} \\
		\hline
		Win\_10\_Client1 & 192.168.1.2         & 255.255.255.0   & 192.168.1.5 (FW1)   & Station       \\
		\hline
		Win\_10\_Client2 & 172.16.3.3          & 255.255.255.0   & 172.16.3.4 (R1)     & Station       \\
		\hline
		Kali\_1          & 172.16.4.4          & 255.255.255.0   & 172.16.4.5 (R1)     & Station       \\
		\hline
		Svr\_Ubuntu      & 192.168.2.2         & 255.255.255.0   & 192.168.2.5 (FW1)   & Serveur       \\
		\hline
	\end{tabular}
\end{table}
\subsection{Présentation de l'infrastructure}
Pour la mise en place et la configuration de notre infrastructure, nous avons
utilisé le logiciel GNS3 qui nous a permis de simuler les équipements réseaux.
L'image ci après présente l'architecture de notre infrastructure réseau.
\clearpage
\begin{figure}[h]
	\centering
	\includegraphics[width=1\textwidth]{Images/Infrastructure.png}
	\caption{Architecture de l'infrastructure réseau}
	\label{fig:infrastructure}
\end{figure}
\section{Recommandations}
Après la configuration de notre infrastructure, nous avons effectué quelques
tests de connectivité entre les machines clientes et le serveur. Nous avons
donc constaté que 04 machines étaient connectées au réseau local (LAN) et 01
machine était connectée à la DMZ. Nous avons donc recommandé les actions
suivantes :
\begin{itemize}
	\item \textbf{Vérification de la connectivité :} S'assurer que toutes les machines
	      sont correctement connectées au réseau et peuvent communiquer entre elles.
	\item \textbf{Sécurisation du réseau :} Mettre en place des règles de sécurité sur le
	      parefeu pour protéger les machines contre les attaques externes.
	\item \textbf{Surveillance régulière:} Vérifier régulièrement l'etat de notre
	      infrastructure pour détecter d'éventuelles anomalies ou pannes de fonctionnement
	\item \textbf{Mises à jour régulières :} Assurer la mise à jour régulière des systèmes
	      d'exploitation de chaque.
\end{itemize}
\end{document}